% GENERATED FILE. Edit canonical lessons, then run npm run generate:books.
% canonical-source-hash: fnv1a64:14ec7e2acec8f7e6
% canonical-lessons: SA-C361-sraddha, SA-C361-nistha, SA-C361-dambhah, SA-C361-krtaghnah, SA-C361-pratisrnoti

\chapter{Faith and Loyalty}
\label{ch:sa-faith-and-loyalty}
\hlchaptermodality{361}

\begin{chapteropening}
\textbf{By the end of this chapter:} \emph{I can talk about faith, steadfastness and ingratitude.}

Complete the last lesson of chapter 361: I can talk about faith, steadfastness and ingratitude.
\end{chapteropening}

\section[\texorpdfstring{\sk{श्रद्धा} (śraddhā)}{श्रद्धा (śraddhā)}]{\sk{श्रद्धा} (śraddhā) — faith; trust}

\label{lesson:SA-C361-sraddha}



\begin{quote}
Before the new one: say the Sanskrit for an insult, then the Sanskrit for misleads.
\end{quote}

\subsection*{You'll want to know: \sk{श्रद्धा}}
\textbf{\sk{श्रद्धा}} — \emph{śraddhā} — \textquotedblleft{}faith; trust\textquotedblright{}.

\begin{quote}
\sk{मम} \sk{वैद्ये} \sk{श्रद्धा} \sk{अस्ति।}
\end{quote}

\begin{quote}
\emph{mama vaidye śraddhā asti.}
\end{quote}

\begin{quote}
I have faith in my doctor.
\end{quote}

A feminine noun.

\begin{cousinweb}
\textbf{\sk{श्रद्धा}} (\emph{śraddhā}) is built from \emph{śrat}, \textquotedblleft{}trust, truth\textquotedblright{}, and the root √dhā, \textquotedblleft{}to put\textquotedblright{}; MW gives \emph{śrad-dhā} \textquotedblleft{}faith, trust, confidence, belief\textquotedblright{} (MW).
\end{cousinweb}

\begin{grammarlens}[title={What you've built}]
One more word for giving reasons and opinions.
\end{grammarlens}

\subsection*{Your turn}
\begin{itemize}
\raggedright
  \item \emph{Say it:} \emph{śraddhā}
  \item \emph{Say it:} \emph{śraddhā}, once more
  \item \emph{Say it:} \emph{mama vaidye śraddhā asti.} — I have faith in my doctor.
  \item \emph{Recall:} say the Sanskrit for an insult, then the Sanskrit for misleads, then say \emph{śraddhā} again
\end{itemize}

\begin{tcolorbox}[breakable,colback=teal!4,colframe=teal!35!black,title={Before you move on}]
What does \sk{श्रद्धा} mean? (\textquotedblleft{}Faith; trust\textquotedblright{}.) Say it once more.
\end{tcolorbox}
\section[\texorpdfstring{\sk{निष्ठा} (niṣṭhā)}{निष्ठा (niṣṭhā)}]{\sk{निष्ठा} (niṣṭhā) — steadfastness; devotion}

\label{lesson:SA-C361-nistha}



\begin{quote}
Before the new one: say the Sanskrit for misleads, then the Sanskrit for faith.
\end{quote}

\subsection*{You'll want to know: \sk{निष्ठा}}
\textbf{\sk{निष्ठा}} — \emph{niṣṭhā} — \textquotedblleft{}steadfastness; devotion\textquotedblright{}.

\begin{quote}
\sk{अरुणस्य} \sk{कार्ये} \sk{निष्ठा} \sk{अस्ति।}
\end{quote}

\begin{quote}
\emph{aruṇasya kārye niṣṭhā asti.}
\end{quote}

\begin{quote}
Arun is devoted to his work.
\end{quote}

A feminine noun.

\begin{cousinweb}
\textbf{\sk{निष्ठा}} (\emph{niṣṭhā}) is built on \emph{ni}, \textquotedblleft{}down\textquotedblright{}, and the root √sthā, \textquotedblleft{}to stand\textquotedblright{}; MW gives \emph{ni-ṣṭhā} \textquotedblleft{}firmness, steadiness, attachment, devotion\textquotedblright{} (MW).
\end{cousinweb}

\begin{grammarlens}[title={What you've built}]
One more word for giving reasons and opinions.
\end{grammarlens}

\subsection*{Your turn}
\begin{itemize}
\raggedright
  \item \emph{Say it:} \emph{niṣṭhā}
  \item \emph{Say it:} \emph{niṣṭhā}, once more
  \item \emph{Say it:} \emph{aruṇasya kārye niṣṭhā asti.} — Arun is devoted to his work.
  \item \emph{Recall:} say the Sanskrit for misleads, then the Sanskrit for faith, then say \emph{niṣṭhā} again
\end{itemize}

\begin{tcolorbox}[breakable,colback=teal!4,colframe=teal!35!black,title={Before you move on}]
What does \sk{निष्ठा} mean? (\textquotedblleft{}Steadfastness; devotion\textquotedblright{}.) Say it once more.
\end{tcolorbox}
\section[\texorpdfstring{\sk{दम्भः} (dambhaḥ)}{दम्भः (dambhaḥ)}]{\sk{दम्भः} (dambhaḥ) — hypocrisy; deceit}

\label{lesson:SA-C361-dambhah}



\begin{quote}
Before the new one: say the Sanskrit for faith, then the Sanskrit for steadfastness.
\end{quote}

\subsection*{You'll want to know: \sk{दम्भः}}
\textbf{\sk{दम्भः}} — \emph{dambhaḥ} — \textquotedblleft{}hypocrisy; deceit\textquotedblright{}.

\begin{quote}
\sk{दम्भः} \sk{सत्यस्य} \sk{शत्रुः} \sk{अस्ति।}
\end{quote}

\begin{quote}
\emph{dambhaḥ satyasya śatruḥ asti.}
\end{quote}

\begin{quote}
Hypocrisy is the enemy of truth.
\end{quote}

A masculine noun.

\begin{cousinweb}
\textbf{\sk{दम्भः}} (\emph{dambhaḥ}) comes from the root √dambh, \textquotedblleft{}to deceive, to injure\textquotedblright{}; MW gives \emph{dambha} \textquotedblleft{}deceit, fraud, feigning, hypocrisy\textquotedblright{} (MW).
\end{cousinweb}

\begin{grammarlens}[title={What you've built}]
One more word for giving reasons and opinions.
\end{grammarlens}

\subsection*{Your turn}
\begin{itemize}
\raggedright
  \item \emph{Say it:} \emph{dambhaḥ}
  \item \emph{Say it:} \emph{dambhaḥ}, once more
  \item \emph{Say it:} \emph{dambhaḥ satyasya śatruḥ asti.} — Hypocrisy is the enemy of truth.
  \item \emph{Recall:} say the Sanskrit for faith, then the Sanskrit for steadfastness, then say \emph{dambhaḥ} again
\end{itemize}

\begin{tcolorbox}[breakable,colback=teal!4,colframe=teal!35!black,title={Before you move on}]
What does \sk{दम्भः} mean? (\textquotedblleft{}Hypocrisy; deceit\textquotedblright{}.) Say it once more.
\end{tcolorbox}
\section[\texorpdfstring{\sk{कृतघ्नः} (kṛtaghnaḥ)}{कृतघ्नः (kṛtaghnaḥ)}]{\sk{कृतघ्नः} (kṛtaghnaḥ) — ungrateful}

\label{lesson:SA-C361-krtaghnah}



\begin{quote}
Before the new one: say the Sanskrit for steadfastness, then the Sanskrit for hypocrisy.
\end{quote}

\subsection*{You'll want to know: \sk{कृतघ्नः}}
\textbf{\sk{कृतघ्नः}} — \emph{kṛtaghnaḥ} — \textquotedblleft{}ungrateful\textquotedblright{}.

\begin{quote}
\sk{कृतघ्नः} \sk{जनः} \sk{उपकारं} \sk{विस्मरति।}
\end{quote}

\begin{quote}
\emph{kṛtaghnaḥ janaḥ upakāraṃ vismarati.}
\end{quote}

\begin{quote}
An ungrateful person forgets a kindness.
\end{quote}

\begin{cousinweb}
\textbf{\sk{कृतघ्नः}} (\emph{kṛtaghnaḥ}) is built from \emph{kṛta}, \textquotedblleft{}what has been done\textquotedblright{}, and the root √han, \textquotedblleft{}to strike, to destroy\textquotedblright{}; MW explains it as \textquotedblleft{}destroying past services\textquotedblright{}, that is \textquotedblleft{}ungrateful\textquotedblright{} (MW).
\end{cousinweb}

\begin{grammarlens}[title={What you've built}]
One more word for giving reasons and opinions.
\end{grammarlens}

\subsection*{Your turn}
\begin{itemize}
\raggedright
  \item \emph{Say it:} \emph{kṛtaghnaḥ}
  \item \emph{Say it:} \emph{kṛtaghnaḥ}, once more
  \item \emph{Say it:} \emph{kṛtaghnaḥ janaḥ upakāraṃ vismarati.} — An ungrateful person forgets a kindness.
  \item \emph{Recall:} say the Sanskrit for steadfastness, then the Sanskrit for hypocrisy, then say \emph{kṛtaghnaḥ} again
\end{itemize}

\begin{tcolorbox}[breakable,colback=teal!4,colframe=teal!35!black,title={Before you move on}]
What does \sk{कृतघ्नः} mean? (\textquotedblleft{}Ungrateful\textquotedblright{}.) Say it once more.
\end{tcolorbox}
\section[\texorpdfstring{\sk{प्रतिशृणोति} (pratiśṛṇoti)}{प्रतिशृणोति (pratiśṛṇoti)}]{\sk{प्रतिशृणोति} (pratiśṛṇoti) — gives one's word; promises}

\label{lesson:SA-C361-pratisrnoti}



\begin{quote}
Before the new one: say the Sanskrit for hypocrisy, then the Sanskrit for ungrateful.
\end{quote}

\subsection*{You'll want to know: \sk{प्रतिशृणोति}}
\textbf{\sk{प्रतिशृणोति}} — \emph{pratiśṛṇoti} — \textquotedblleft{}gives one's word; promises\textquotedblright{}.

\begin{quote}
\sk{मीरा} \sk{साहाय्यं} \sk{प्रतिशृणोति।}
\end{quote}

\begin{quote}
\emph{mīrā sāhāyyaṃ pratiśṛṇoti.}
\end{quote}

\begin{quote}
Mira gives her word that she will help.
\end{quote}

Past (he …): \sk{सः} \sk{साहाय्यं} \sk{प्रतिश्रुतवान्}, \emph{saḥ sāhāyyaṃ pratiśrutavān}.

\begin{cousinweb}
\textbf{\sk{प्रतिशृणोति}} (\emph{pratiśṛṇoti}) is built on \emph{prati}, \textquotedblleft{}towards\textquotedblright{}, and the root √śru, \textquotedblleft{}to hear\textquotedblright{}; MW gives \emph{prati-√śru} \textquotedblleft{}to assure, to agree, to promise\textquotedblright{} (MW).
\end{cousinweb}

\begin{grammarlens}[title={What you've built}]
One more word for giving reasons and opinions.
\end{grammarlens}

\subsection*{Your turn}
\begin{itemize}
\raggedright
  \item \emph{Say it:} \emph{pratiśṛṇoti}
  \item \emph{Say it:} \emph{pratiśṛṇoti}, once more
  \item \emph{Say it:} \emph{mīrā sāhāyyaṃ pratiśṛṇoti.} — Mira gives her word that she will help.
  \item \emph{Recall:} say the Sanskrit for hypocrisy, then the Sanskrit for ungrateful, then say \emph{pratiśṛṇoti} again
\end{itemize}

\begin{tcolorbox}[breakable,colback=teal!4,colframe=teal!35!black,title={Before you move on}]
What does \sk{प्रतिशृणोति} mean? (\textquotedblleft{}Gives one's word; promises\textquotedblright{}.) Say it once more.
\end{tcolorbox}
